% Corte mecánico íntegro: parte_iv.tex:107--219.
% Residencia material del ensamblaje sucesor de 102 capítulos.
% Residencia de realización: capítulo 63.
\chapter{Réalisation électromagnétique restreinte}
\label{espiral:chap:realizacion-em}

\section{Publication dimensionnelle ultérieure}

Après construction du continu et dérivation de l'opérateur constitutif du vide, la vitesse de propagation est publiée sous la forme
\[
 c_{\mathrm{vac}}=(\mu_{\mathrm{vac}}\varepsilon_{\mathrm{vac}})^{-1/2}.
\]
Soit
\[
 \zeta=\frac{z-c_{\mathrm{vac}}t}{\lambda}
\]
et \(\sigma=\pm1\). Definimos
\begin{align}
 \mathbf E_\sigma(z,t)
 &=E_0\bigl(
 \mathbf e_1\cos2\pi\zeta
 +\sigma\mathbf e_2\sin2\pi\zeta\bigr),\\
 \mathbf B_\sigma(z,t)
 &=\frac1{c_{\mathrm{vac}}}\widehat{\mathbf k}
 \times\mathbf E_\sigma(z,t).
 \label{espiral:eq:modo-em-circular}
\end{align}
Alors
\[
 \mathbf E_\sigma\perp\mathbf B_\sigma\perp\widehat{\mathbf k},
 \qquad
 |\mathbf E_\sigma|=c_{\mathrm{vac}}|\mathbf B_\sigma|,
\]
\[
 \omega=c_{\mathrm{vac}}k,
 \qquad
 \mathbf S=\frac1{\mu_{\mathrm{vac}}}
 \mathbf E_\sigma\times\mathbf B_\sigma
 =\frac{E_0^2}{Z_{\mathrm{vac}}}\widehat{\mathbf k}.
\]

\section{Réalisation hélicoïdale de la phase électromagnétique : retour périodique et translation longitudinale d'une longueur d'onde}

À temps fixé, l'extrémité normalisée du vecteur électrique satisfait
\begin{equation}
 \left(
 \frac{E_1}{E_0},\frac{\sigma E_2}{E_0},
 \frac{z-c_{\mathrm{vac}}t}{\lambda}
 \right)
 =\bigl(\cos2\pi\zeta,\sin2\pi\zeta,\zeta\bigr).
 \label{espiral:eq:helice-em-fase-posicion}
\end{equation}
L'équation coïncide avec l'hélice normalisée de phase et de mémoire. L'étape \(\zeta\mapsto\zeta+1\) ramène la phase électromagnétique et translate l'événement marqué d'une longueur d'onde.

\begin{figure}[htbp]
\centering
\begin{tikzpicture}[x={(1cm,0cm)},y={(.28cm,.12cm)},z={(0cm,.65cm)},scale=.77,>=Latex]
 \draw[->,HMTGray,thick] (0,0,0)--(0,0,6.8) node[above] {$\zeta$};
 \draw[HMTBlue,very thick,domain=0:1080,samples=220,smooth,variable=\t]
   plot ({1.35*cos(\t)},{1.35*sin(\t)},{\t/170});
 \foreach \k in {0,180,...,900}{
   \draw[->,HMTRed,thick] (0,0,{\k/170})--({1.35*cos(\k)},{1.35*sin(\k)},{\k/170});
   \draw[->,HMTGold,thick] (0,0,{\k/170})--({-1.0*sin(\k)},{1.0*cos(\k)},{\k/170});
 }
 \node[HMTRed,font=\small\sffamily] at (1.8,0,5.5) {$\mathbf E$};
 \node[HMTGold,font=\small\sffamily] at (-1.4,0,4.5) {$\mathbf B$};
\end{tikzpicture}
\caption{Réalisation analytique de l'hélice phase--position dans un mode circulaire.}
\label{espiral:fig:helice-em}
\end{figure}

L'affirmation exacte porte sur la trajectoire de l'extrémité du champ dans le fibré phase--position. On ne conclut pas que toute ligne spatiale de champ est une hélice. On n'identifie pas non plus l'hélicité électromagnétique à l'holonomie fermionique de l'électron.

\section{Dualité électromagnétique}

Sur les deux-formes complexifiées en signature lorentzienne,
\[
 \star^2=-I.
\]
La relation
\[
 \mathfrak R_{\mathrm{EM}}\circ J_{+1}
 =\star\circ\mathfrak R_{\mathrm{EM}}
\]
réalise le quart de tour interne comme dualité électrique--magnétique lorsque l'opérateur d'entrelacement \(\mathfrak R_{\mathrm{EM}}\) est construit. Les espaces propres
\[
 \star F_\sigma=\ii\sigma F_\sigma
\]
séparent les deux hélicités. C'est une interface physique typée : elle utilise le générateur elliptique déjà construit, mais requiert l'espace des formes et la métrique de la réalisation.

\section{Condition de compatibilité avec la vis de Pell}

Pour une phase
\[
 \Phi(r,\varphi,z,t)
 =\ell\varphi+\nu\log(r/r_0)+k_zz-\omega t,
\]
l'invariance sous \((r,\varphi)\mapsto(\varepsilon r,\varphi+\pi/2)\) exige
\begin{equation}
 \nu\log(2+\sqrt3)+\ell\frac\pi2
 \in2\pi\Z.
 \label{espiral:eq:cuantizacion-tornillo-em}
\end{equation}
Cette condition quantifie la compatibilité de phase. L'amplitude vectorielle doit satisfaire en outre l'équation de Helmholtz, la transversalité et les conditions au bord.
